% R. Brown's translations of two letters to L. Breen, with his note and comments (cote 134-2): archive cote 134-2, site page 125 (PDF page 126).
% translation of the letter of 17/2/1975. Draft compared: none.
% AI-transcribed by Claude Opus 5.5 (claude-opus-5-5, adaptive thinking, effort high), Message Batches API, batch brown_full (msgbatch_01Md3i5yK1PxkpJn5GHfEExE), 2026-11-13; not yet checked by eye.
\scanpage{134-2/125}
\mspage{- 4 -}

same way that this induces isomorphisms on $\pi_i$ for $1 \leq i \leq n + 1$ .]

[\uline{N.B}. One sees by the way that for $F$ an n-champ of groupoids on $X$ ,
if one restricts to defining the $H^i(X, F)$ for $0 \leq i \leq n$ , one has no need
of a Picard structure on $F$ , as it is sufficient to put
\[ H^i(X, F) = \pi_{n-i}(\Gamma_X(F)) \qquad 0 \leq i \leq n \ . \]
If on the other hand $F$ is an n-Gr-champ (i.e. $F$ has the structure of a
composition law $F \times F \rightarrow F$ with the usual formal properties of a group) the
``classifying (n+1)-champ'' is defined, and one can define $H^i(X, F)$ for
$0 \leq i \leq n + 1$ by
\[ H^i(X, F) = \pi_{n+1-i}(\Gamma_X(\mathbb{C}F)) \]
in particular
\[ H^{n+1}(X, F) = \pi_0(\Gamma_X(\mathbb{C}F)) \]
\[ \hspace{6em} = \text{equivalence classes of sections of } \mathbb{C}F \ . \]
But one can form $C\mathbb{C}F = \mathbb{C}^2F$ and define $H^{n+2}(X, F)$ , it seems \textit{only} if $\mathbb{C}F$
is itself a Gr-(n+1)-champ, which is without doubt the case only if $F$ is
a strict Picard n-champ ....

Let us now come to the case where $F$ is a \textit{locally constant} n-champ on
$X$ , and so is defined by an (n+1)-functor
\[ C_{n+1} \xrightarrow{F} \text{(strict Picard n-Cat)} \ . \]

Then, putting for $0 \leq i \leq n$
\[ H^i(C_{n+1}, F) = \pi_{n-i}(\mathit{Hom}(e_n^{(C_{n+1})}, F)) \ , \]
``one knows it fails'', as one has a canonical isomorphism
\[ H^i(C_{n+1}, F) \cong H^i(X, F) \]

valid in effect without Picard structure on $F$ ... It is thus necessary
for all $i$ and for every $\infty$-groupoid $C$ and every (n+1)-functor
\[ C \xrightarrow{F} \text{(strict Picard n-Cat)} \ , \]
to define
\[ H^i(C, F) = \pi_{N-i}\mathit{Hom}(e_N^C, \mathbb{C}^{N-n}F) \]

where one chooses $N \geq \sup(i, n)$ , If $F$ has only a Gr-structure (not
necessarily Picard) one can define the $H^i(C, F)$ for $i \leq n + 1$ by
\[ H^i(C, F) = \pi_{n+1-i}\mathit{Hom}(e^C_{n+1}, \mathbb{C}F) \ . \]
In the case $C = C_{n+1} = \Pi_{n+1}(X)$ , it must still be true (by virtue of (A)
above), that this set is canonically isomorphic to $H^{n+1}(X, F) = \pi_0\Gamma_{\unsure{x|\times}}(\mathbb{C}F)$
(this is true and very easy for $n = 0$). Can one describe the arrow
between the two sides of
\[ H^{n+1}(X, F) \simeq H^{n+1}(\Pi_{n+1}(X), F) \ ? \]

If one wishes to make (A) and (B) explicit again, in terms of the yoga
(C), one comes to the following situation:

\mspage{\hand{126}}
